\section{Метод конечных элементов. Эрмитовы элементы}
\label{sec:fem-hermite}

\subsection{Назначение эрмитовых элементов}

Лагранжевы элементы обеспечивают непрерывность функции, но производные между
элементами обычно разрывны. Для задач, слабая постановка которых требует
пространства $H^2$, нужна $C^1$-согласованная аппроксимация. Типичный пример ---
уравнение изгиба балки Эйлера--Бернулли четвёртого порядка:
\[
    \bigl(EIu''\bigr)''=q(x).
\]

Эрмитовы элементы используют в качестве степеней свободы не только значения
функции, но и значения её производных в узлах. Это позволяет обеспечить
непрерывность первой производной.

\subsection{Слабая постановка задачи изгиба балки}

Для простоты рассмотрим
\[
    EIu^{(4)}=q(x),
    \qquad x\in(0,L).
\]
После двукратного интегрирования по частям получаем билинейную форму
\[
    a(u,v)=\int_0^L EIu''v''\,dx.
\]
Слабая постановка требует
\[
    u,v\in H^2(0,L)
    \quad\text{с учётом существенных граничных условий}.
\]
Следовательно, конечно-элементное пространство должно быть глобально
$C^1$-непрерывным.

\subsection{Кубический эрмитов элемент в одном измерении}

На элементе $K=[x_1,x_2]$ длины $h$ используются четыре степени свободы:
\[
    u(x_1),\qquad u'(x_1),\qquad u(x_2),\qquad u'(x_2).
\]
В локальной координате
\[
    \xi=\frac{x-x_1}{h}\in[0,1]
\]
приближение записывается как
\[
    u_h(x)=
    N_1(\xi)u_1+N_2(\xi)\theta_1+
    N_3(\xi)u_2+N_4(\xi)\theta_2,
\]
где $\theta_i=u'(x_i)$, а функции формы равны
\[
\begin{aligned}
N_1(\xi)&=1-3\xi^2+2\xi^3,\\
N_2(\xi)&=h(\xi-2\xi^2+\xi^3),\\
N_3(\xi)&=3\xi^2-2\xi^3,\\
N_4(\xi)&=h(-\xi^2+\xi^3).
\end{aligned}
\]
Они удовлетворяют интерполяционным условиям:
\[
\begin{array}{c|cccc}
&N_1&N_2&N_3&N_4\\ \hline
x=x_1&1&0&0&0\\
\dfrac{d}{dx}\big|_{x_1}&0&1&0&0\\
x=x_2&0&0&1&0\\
\dfrac{d}{dx}\big|_{x_2}&0&0&0&1
\end{array}
\]

\subsection{Локальная матрица жёсткости балки}

При постоянном $EI$ локальная матрица для вектора степеней свободы
\[
    U_e=(u_1,\theta_1,u_2,\theta_2)^T
\]
имеет вид
\[
K^{(e)}=\frac{EI}{h^3}
\begin{pmatrix}
12&6h&-12&6h\\
6h&4h^2&-6h&2h^2\\
-12&-6h&12&-6h\\
6h&2h^2&-6h&4h^2
\end{pmatrix}.
\]
Она получается из формулы
\[
    K_{ij}^{(e)}=\int_{x_1}^{x_2}EI N_i''N_j''\,dx.
\]
Локальный вектор нагрузки:
\[
    F_i^{(e)}=\int_{x_1}^{x_2}q(x)N_i(x)\,dx.
\]
После сборки значения перемещения и угла поворота в общем узле совпадают для
соседних элементов, поэтому глобальная функция принадлежит $C^1$.

\subsection{Существенные и естественные условия}

Для балки геометрические условия
\[
    u=\overline u,
    \qquad
    u'=\overline\theta
\]
являются существенными и задаются на степенях свободы. Момент и поперечная сила,
связанные с
\[
    M=EIu'',
    \qquad
    Q=-(EIu'')',
\]
возникают как естественные граничные условия после интегрирования по частям.

\subsection{Эрмитова интерполяция и порядок точности}

Кубический эрмитов интерполянт точно воспроизводит полиномы степени не выше
трёх. Для достаточно гладкой функции локальная ошибка имеет порядок
\[
    \|u-I_hu\|_{L^{\infty}(K)}=O(h^4),
\]
а ошибка производной --- $O(h^3)$. В энергетической норме задачи балки при
регулярной сетке получают типичную оценку порядка $O(h^2)$.

\subsection[Многомерные C1-элементы]{Многомерные $C^1$-элементы}

В двух измерениях построение $C^1$-согласованных элементов значительно сложнее.
Примеры:
\begin{itemize}
    \item прямоугольный элемент Богнера--Фокса--Шмита;
    \item треугольный элемент Аргириса пятой степени;
    \item элемент Клафа--Точера с разбиением треугольника.
\end{itemize}
Они используются для пластин, оболочек и уравнений четвёртого порядка.
На практике часто применяют смешанные постановки, позволяющие снизить порядок
уравнения и использовать обычные $C^0$-элементы.

\subsection{Сравнение с лагранжевыми элементами}

Лагранжевы элементы:
\begin{itemize}
    \item степени свободы --- значения функции;
    \item глобальная непрерывность --- обычно $C^0$;
    \item подходят для уравнений второго порядка.
\end{itemize}

Эрмитовы элементы:
\begin{itemize}
    \item степени свободы --- значения функции и производных;
    \item обеспечивают $C^1$-непрерывность;
    \item естественны для задач четвёртого порядка, но сложнее в многомерном случае.
\end{itemize}

\subsection{Полная вариационная формула для балки}

Для балки Эйлера--Бернулли
\[
    (EIu'')''=q
\]
умножим уравнение на $v$ и дважды проинтегрируем по частям:
\[
 \int_0^L EIu''v''\,dx
 =\int_0^L qv\,dx
 +\left[(EIy'')'v-EIy''v'\right]_{0}^{L}.
\]
В зависимости от закрепления существенными условиями являются значения прогиба $u$ и
угла поворота $u'$, а естественными --- изгибающий момент
\[
    M=EIu''
\]
и поперечная сила
\[
    Q=-(EIu'')'
\]
с учётом принятого соглашения знаков. Поэтому пространство $H^2$ возникает не
искусственно: слабая форма содержит вторые производные решения и тестовой функции.

\subsection{Эрмитова интерполяция как условие унисольвентности}

Кубический многочлен на элементе однозначно определяется четырьмя степенями свободы
\[
    u(x_1),\quad u'(x_1),\quad u(x_2),\quad u'(x_2).
\]
Это пример \textbf{унисольвентного} набора степеней свободы: если все четыре функционала
обращаются в нуль на кубическом многочлене, то сам многочлен тождественно равен нулю.
Отсюда следует существование единственного набора функций формы, двойственного этим
степеням свободы.

На соседних элементах глобальные значения $u$ и $u'$ отождествляются. Именно это
обеспечивает $C^1$-непрерывность и принадлежность конечно-элементной функции пространству
$H^2$ в одномерной задаче.

\subsection{Порядок сходимости для балки}

Кубическая эрмитова интерполяция для достаточно гладкой функции имеет типичные оценки
\[
    \|u-I_hu\|_{L^2}=O(h^4),\qquad
    \|u'- (I_hu)'\|_{L^2}=O(h^3),\qquad
    \|u''-(I_hu)''\|_{L^2}=O(h^2).
\]
Энергетическая норма балки определяется второй производной, поэтому естественный
порядок ошибки энергии для кубических эрмитовых элементов составляет $O(h^2)$ при
достаточной гладкости решения.

Для задач второго порядка такие элементы обычно избыточны; их преимущество проявляется
именно в задачах четвёртого порядка и в моделях, где требуется согласованная
аппроксимация производных.

\subsection{Что важно сказать на экзамене}

Главное отличие эрмитовых элементов --- наличие производных среди степеней
свободы. На примере кубического элемента балки нужно привести четыре функции
формы, объяснить $C^1$-согласованность и показать, почему она требуется слабой
постановкой уравнения четвёртого порядка.

\newpage